\documentclass[11pt]{article}
\usepackage[margin=25mm]{geometry}
\usepackage{amsmath,amssymb,amsthm}
\usepackage[T1]{fontenc}
\usepackage[colorlinks=true,urlcolor=blue]{hyperref}
\newtheorem{theorem}{Theorem}
\setlength{\parindent}{0pt}\setlength{\parskip}{5pt}
\begin{document}
\begin{center}{\Large Equal expressive power preserves convergence laws}\\[5pt]{\large Conjecture 00000005030}\\[5pt]zbw123456 \quad 14 October 2026\end{center}

Fix a common class $\Omega$ of structures and a common sequence $(\mu_n)$ of probability distributions. A sentence $\varphi$ defines its truth event $[\varphi]\subseteq\Omega$. Structures of varying finite sizes may be placed in their disjoint union, with $\mu_n$ supported on the $n$-th part.

Two logics have the same expressive power when every sentence in either logic has an equivalent sentence in the other, uniformly on all structures in this class. A convergence law means that $\mu_n([\varphi])$ converges for every fixed sentence.
\begin{theorem}
Logics of the same expressive power satisfy a convergence law simultaneously on the same random structures.
\end{theorem}
\begin{proof}
Suppose $L$ satisfies the convergence law and let $\psi$ be a sentence of $K$. Equal expressive power gives a sentence $\varphi$ of $L$ with $[\psi]=[\varphi]$. Therefore
\[
\mu_n([\psi])=\mu_n([\varphi])\quad\text{for every }n.
\]
The right-hand sequence converges, so the left-hand sequence does. This proves the convergence law for $K$. Interchanging the two logics gives the reverse implication. Thus no pair of counting extensions, or any other pair of logics, can have the separation claimed in the problem under these definitions.
\end{proof}
\textbf{Meaning of equality.} The equivalent sentence is fixed independently of $n$. Equivalence only at each fixed size with size-dependent translations is not equality of expressive power. Changing the probability distributions also changes the random-structure model. Neither alteration occurs in the conjecture.

\textbf{Lean correspondence.} \texttt{Logic} stores a sentence type, satisfaction predicate and measurable truth events. \texttt{AtMostExpressive} is mutual definability in one direction; \texttt{SameExpressivePower} contains both directions. \texttt{ConvergenceLaw} uses Mathlib probability measures, real probabilities and filter convergence. \texttt{transfer} proves the one-way transfer by equality of events; \texttt{no\_separation} and \texttt{conjecture\_false} give the equivalence and negate the asserted existential separation. No convergence theorem about a particular distribution is assumed.


\textbf{Verification.} Lean 4.19.0 with Mathlib pinned to
\texttt{c44e0c8ee63ca166450922a373c7409c5d26b00b}.
The accompanying project and verification transcript give the exact theorem and axiom audit. No proof placeholders, native evaluator proof bypass, or custom axiom declarations are used. This is a prepared submission; official acceptance has not been claimed.
\end{document}
